% ============================================================================
% 08_diagramas_colaboracion.tex — Cuerpo del entregable "Diagramas de
% Colaboración". Sin preámbulo: se incluye tanto desde
% 08_Diagramas_Colaboracion.tex (standalone) como, eventualmente, desde un
% documento maestro.
% ============================================================================
\section{Diagramas de Colaboración}

\subsection{Propósito y alcance}

Esta sección documenta, siguiendo la técnica descripta por Craig Larman en
\textit{UML y Patrones}, los \textbf{contratos de operación de sistema} y los
\textbf{diagramas de colaboración} (\textit{communication diagrams} en UML
2.5) correspondientes a las ocho operaciones de sistema identificadas en
\texttt{docs/IDS.md}: SO-01 a SO-05 (implementadas) y SO-06 a SO-08
(\textbf{diseñadas, pendientes de implementación}, cubren RF-03, RF-04 y
RF-06 oficiales). Cada operación de sistema es un mensaje
que un actor le envía al sistema considerado como una caja negra durante el
diseño de la interacción; el diagrama de colaboración muestra cómo el
sistema, ya descompuesto en sus objetos de diseño reales (\textit{hooks} de
React, componentes, \textit{handlers} de FastAPI y el modelo de
reconocimiento), coopera para cumplir esa operación.

A diferencia de un diagrama de secuencia, un diagrama de colaboración no
tiene eje temporal ni líneas de vida: prioriza la topología de
colaboración entre objetos. Cada par de objetos que colabora está unido por
un único enlace estructural, dibujado como una línea recta \textbf{sin
punta de flecha}. Cada mensaje que viaja por ese enlace se dibuja como una
flechita corta y propia, paralela al enlace, apuntando siempre hacia el
objeto receptor, con su número jerárquico junto al texto del mensaje (1,
1.1, 1.2, 2, 2.1\ldots); cuando varios mensajes recorren el mismo enlace se
apilan sobre él, y los mensajes de sentido opuesto (por ejemplo, un mensaje
de vuelta genuino, como el evento asíncrono de un WebSocket) se dibujan del
lado contrario del enlace. Los valores de retorno de una llamada síncrona
simple no se dibujan como una flecha aparte: se integran en el propio
mensaje con la notación \texttt{resultado := operación(args)}, siguiendo la
convención de Larman de omitir flechas de retorno triviales. Un
automensaje (un objeto que se invoca a sí mismo) se dibuja como un pequeño
bucle sobre el propio objeto. Cada objeto se identifica como
\texttt{instancia~:~Clase} (o \texttt{:Clase} cuando no hace falta nombrar
la instancia), con el nombre subrayado, cuando corresponde a una instancia
real (por ejemplo, un componente de React o el modelo LSTM cargado en
memoria), o simplemente por su nombre cuando representa un módulo o una
colección de funciones del backend sin una clase propia (por ejemplo,
\texttt{capture\_validation} o \texttt{MEDICAL\_TEMPLATES}), ya que el
backend de LSA Salud (\texttt{app.py}) no está organizado en clases salvo por el
modelo \texttt{SignLanguageLSTM}. Los cinco diagramas se dibujan en
draw.io (formato \texttt{.drawio} editable, generado desde
\texttt{docs/diagramas/colaboracion/gen\_drawio.py}) para poder expresar
esta notación de forma exacta; ver \texttt{docs/README.md} para el flujo de
edición y reexportación.

Para cada operación se presenta: el contrato de la operación (con sus
referencias cruzadas a caso de uso y requisito funcional), el diagrama de
colaboración, los patrones GRASP aplicados y una tabla que traza cada objeto
del diagrama a su archivo fuente real.

\subsection{SO-01: reconocerSeña(frame) --- DC-01}

\subsubsection{Contrato de operación}

\begin{table}[H]
\centering
\begin{tabularx}{\textwidth}{@{}p{4.4cm} X@{}}
\toprule
\textbf{Operación} & \texttt{reconocerSeña(frame)} \\
\textbf{Referencias cruzadas} & Realiza CU-01 (pasos de captura, reconocimiento
y armado de la secuencia); RF-01.1, RF-01.2, RF-01.3, RF-01.4. \\
\textbf{Precondiciones} & La cámara está autorizada y el \textit{hook} de
MediaPipe está emitiendo \textit{frames}; el backend tiene el modelo
\texttt{SignLanguageLSTM} cargado; existe una conexión WebSocket activa hacia
\texttt{/ws/predict}. \\
\textbf{Postcondiciones} & El \textit{frame} quedó agregado al búfer
deslizante de la Estación del Paciente. Si el búfer alcanzó el múltiplo de
\texttt{PREDICT\_STRIDE\_FRAMES} configurado, se envió una predicción al
backend y, si la seña reconocida superó el umbral de confianza y no fue
``idle'', la palabra reconocida quedó agregada a la secuencia en curso
(\texttt{transcript}) de \texttt{useSentenceBuilder}. \\
\bottomrule
\end{tabularx}
\caption{Contrato de la operación SO-01 --- reconocerSeña(frame)}
\label{tab:col-contrato-so01}
\end{table}

\subsubsection{Diagrama de colaboración DC-01}

\diagrama{08_colaboracion_dc01}{\textwidth}{DC-01 --- Colaboración para reconocerSeña(frame)}{fig:col-dc01}

\subsubsection{Patrones GRASP aplicados}

\begin{itemize}[leftmargin=*, itemsep=2pt]
    \item \textbf{Controlador:} \texttt{handleFrame} en \texttt{PacientePage}
    actúa como controlador de la operación de sistema: recibe el evento de
    captura desde el \textit{hook} de MediaPipe y coordina al búfer, al
    constructor de oraciones y al \textit{hook} de WebSocket, sin contener
    él mismo la lógica de reconocimiento.
    \item \textbf{Experto en información:} el modelo \texttt{SignLanguageLSTM}
    es quien tiene el conocimiento (los pesos entrenados) para clasificar el
    tensor de entrada, por lo que es el responsable de producir la
    predicción; ningún otro objeto duplica esa responsabilidad.
    \item \textbf{Bajo acoplamiento:} \texttt{useSignWebSocket} aísla a
    \texttt{PacientePage} del protocolo de transporte (JSON o binario) y de
    los detalles de reconexión del WebSocket; \texttt{PacientePage} solo
    conoce \texttt{sendFrames} y un \textit{callback} \texttt{onMessage}.
    \item \textbf{Indirección:} el \textit{handler}
    \texttt{websocket\_predict} media entre el frontend y el modelo
    \texttt{SignLanguageLSTM}: el frontend nunca invoca al modelo
    directamente, sino a través de este intermediario, lo que permite
    cambiar el modelo o el preprocesamiento sin afectar al cliente.
\end{itemize}

\subsubsection{Objetos del diagrama y su ubicación en el código}

\begin{table}[H]
\centering
\begin{tabularx}{\textwidth}{@{}l X@{}}
\toprule
\textbf{Objeto} & \textbf{Ubicación real} \\
\midrule
\texttt{useMediaPipe} & Hook (pipeline por defecto) $\rightarrow$ \texttt{web/hooks/useMediaPipe.ts}; intercambiable por \texttt{useMediaPipeTasks} ($\rightarrow$ \texttt{web/hooks/useMediaPipeTasks.ts}) activando \texttt{?mp=tasks} en \texttt{web/app/paciente/page.tsx} --- ambos exponen el mismo contrato \texttt{onFrame} \\
\texttt{handleFrame : PacientePage} & Componente $\rightarrow$ \texttt{web/app/paciente/page.tsx} \\
\texttt{recordedFrames} & Ref (búfer) $\rightarrow$ \texttt{web/app/paciente/page.tsx} \\
\texttt{useSignWebSocket} & Hook $\rightarrow$ \texttt{web/hooks/useSignWebSocket.ts} \\
\texttt{websocket\_predict} & Handler FastAPI $\rightarrow$ \texttt{app.py} \\
\texttt{model : SignLanguageLSTM} & Instancia del modelo $\rightarrow$ \texttt{src/model/lstm.py} \\
\texttt{TopPredictions} & Componente $\rightarrow$ \texttt{web/components/*/TopPredictions.tsx} \\
\texttt{useSentenceBuilder} & Hook $\rightarrow$ \texttt{web/hooks/useSentenceBuilder.ts} \\
\bottomrule
\end{tabularx}
\caption{Trazabilidad de objetos --- DC-01}
\label{tab:col-objetos-dc01}
\end{table}

\subsection{SO-02: cerrarOracionYTraducir(secuencia) --- DC-02}

\subsubsection{Contrato de operación}

\begin{table}[H]
\centering
\begin{tabularx}{\textwidth}{@{}p{4.4cm} X@{}}
\toprule
\textbf{Operación} & \texttt{cerrarOracionYTraducir(secuencia)} \\
\textbf{Referencias cruzadas} & Realiza CU-01 (cierre de secuencia) y CU-02
completo; RF-01.4, RF-01.5. \\
\textbf{Precondiciones} & La secuencia de palabras reconocidas
(\texttt{transcript}) no está vacía; transcurrieron al menos 800~ms de
inactividad y la magnitud de movimiento promedio está por debajo del umbral
de quietud. \\
\textbf{Postcondiciones} & La secuencia en curso quedó cerrada y vaciada
(\texttt{transcriptRef} vuelve a estar vacío). Se obtuvo una oración natural
en español --- por plantilla médica, por caché de traducciones o por el
modelo de lenguaje local --- y, si se obtuvo una oración válida, el mensaje
del paciente quedó agregado a la conversación compartida a través de
\texttt{useChatChannel}. \\
\bottomrule
\end{tabularx}
\caption{Contrato de la operación SO-02 --- cerrarOracionYTraducir(secuencia)}
\label{tab:col-contrato-so02}
\end{table}

\subsubsection{Diagrama de colaboración DC-02}

\diagrama{08_colaboracion_dc02}{\textwidth}{DC-02 --- Colaboración para cerrarOracionYTraducir(secuencia)}{fig:col-dc02}

\subsubsection{Patrones GRASP aplicados}

\begin{itemize}[leftmargin=*, itemsep=2pt]
    \item \textbf{Experto en información:} \texttt{useSentenceBuilder} decide
    cuándo cerrar la oración porque es quien tiene el estado de quietud
    (ventana de magnitudes de movimiento) y el \textit{transcript} en curso;
    ningún otro objeto podría tomar esa decisión sin duplicar ese estado.
    \item \textbf{Controlador:} \texttt{process\_sentence} en \texttt{app.py}
    coordina la resolución de la traducción (plantilla, caché, modelo de
    lenguaje) como punto único de entrada del backend para esta operación.
    \item \textbf{Indirección:} \texttt{MEDICAL\_TEMPLATES} y
    \texttt{\_translation\_cache} median entre la secuencia de palabras y el
    modelo de lenguaje externo, evitando invocar a Ollama cuando ya existe
    una traducción conocida o reciente.
    \item \textbf{Bajo acoplamiento:} \texttt{useSentenceBuilder} no conoce
    cómo se traduce la secuencia (plantilla, caché o LLM); solo conoce el
    contrato HTTP de \texttt{/process\_sentence}, lo que permite cambiar la
    estrategia de traducción en el backend sin tocar el frontend.
\end{itemize}

\subsubsection{Objetos del diagrama y su ubicación en el código}

\begin{table}[H]
\centering
\begin{tabularx}{\textwidth}{@{}l X@{}}
\toprule
\textbf{Objeto} & \textbf{Ubicación real} \\
\midrule
\texttt{useSentenceBuilder} & Hook $\rightarrow$ \texttt{web/hooks/useSentenceBuilder.ts} \\
\texttt{process\_sentence} & Handler FastAPI $\rightarrow$ \texttt{app.py} \\
\texttt{MEDICAL\_TEMPLATES} & Diccionario de módulo $\rightarrow$ \texttt{app.py} \\
\texttt{\_translation\_cache} & Caché LRU en memoria $\rightarrow$ \texttt{app.py} \\
\texttt{Ollama (qwen2:1.5b)} & Proceso local externo $\rightarrow$ API HTTP en \texttt{localhost:11434} \\
\texttt{useChatChannel (Paciente)} & Hook $\rightarrow$ \texttt{web/hooks/useChatChannel.ts} \\
\bottomrule
\end{tabularx}
\caption{Trazabilidad de objetos --- DC-02}
\label{tab:col-objetos-dc02}
\end{table}

\subsection{SO-03: sintetizarVoz(oracion) --- DC-03}

\subsubsection{Contrato de operación}

\begin{table}[H]
\centering
\begin{tabularx}{\textwidth}{@{}p{4.4cm} X@{}}
\toprule
\textbf{Operación} & \texttt{sintetizarVoz(oracion)} \\
\textbf{Referencias cruzadas} & Realiza CU-03 completo; RF-01.6. \\
\textbf{Precondiciones} & Existe una oración traducida no vacía
(resultado de SO-02); el motor Piper TTS (\texttt{piper\_voice}) puede estar
cargado o no. \\
\textbf{Postcondiciones} & Si el motor de síntesis está disponible, el audio
de la oración quedó disponible en formato \textit{base64} --- reutilizado
desde \texttt{\_audio\_cache} si ya existía, o generado y agregado a la
caché en caso contrario --- y se programó su reproducción en el navegador
tras la ventana de exhibición previa. Si el motor no está disponible, el
mensaje queda entregado sin audio, sin interrumpir el flujo de la
conversación. \\
\bottomrule
\end{tabularx}
\caption{Contrato de la operación SO-03 --- sintetizarVoz(oracion)}
\label{tab:col-contrato-so03}
\end{table}

\subsubsection{Diagrama de colaboración DC-03}

\diagrama{08_colaboracion_dc03}{\textwidth}{DC-03 --- Colaboración para sintetizarVoz(oracion)}{fig:col-dc03}

\subsubsection{Patrones GRASP aplicados}

\begin{itemize}[leftmargin=*, itemsep=2pt]
    \item \textbf{Experto en información:} \texttt{piper\_voice} es el único
    objeto con el conocimiento (modelo de voz \texttt{es\_AR-daniela-high})
    necesario para sintetizar el audio, por lo que es el responsable de
    generarlo.
    \item \textbf{Indirección:} \texttt{\_audio\_cache} media entre
    \texttt{process\_sentence} y \texttt{piper\_voice}, evitando resintetizar
    audio ya generado para la misma oración.
    \item \textbf{Bajo acoplamiento:} \texttt{useSentenceBuilder} no conoce
    el motor de síntesis ni el formato interno del audio; solo recibe
    \texttt{audio\_base64} y decide cuándo reproducirlo, tras la ventana de
    exhibición configurada por \texttt{DISPLAY\_WINDOW\_MS}.
    \item \textbf{Alta cohesión:} toda la responsabilidad de obtener el
    audio (consultar caché, sintetizar, cachear) queda concentrada en
    \texttt{process\_sentence}, sin dispersarse hacia el frontend.
\end{itemize}

\subsubsection{Objetos del diagrama y su ubicación en el código}

\begin{table}[H]
\centering
\begin{tabularx}{\textwidth}{@{}l X@{}}
\toprule
\textbf{Objeto} & \textbf{Ubicación real} \\
\midrule
\texttt{process\_sentence} & Handler FastAPI $\rightarrow$ \texttt{app.py} \\
\texttt{\_audio\_cache} & Caché LRU en memoria (hasta 50 oraciones) $\rightarrow$ \texttt{app.py} \\
\texttt{piper\_voice : PiperVoice} & Instancia del motor TTS $\rightarrow$ \texttt{app.py} \\
\texttt{useSentenceBuilder} & Hook $\rightarrow$ \texttt{web/hooks/useSentenceBuilder.ts} \\
\texttt{Audio} & API nativa del navegador, invocada desde \texttt{web/hooks/useSentenceBuilder.ts} \\
\bottomrule
\end{tabularx}
\caption{Trazabilidad de objetos --- DC-03}
\label{tab:col-objetos-dc03}
\end{table}

\subsection{SO-04: enviarRespuesta(mensaje) --- DC-04}

\subsubsection{Contrato de operación}

\begin{table}[H]
\centering
\begin{tabularx}{\textwidth}{@{}p{4.4cm} X@{}}
\toprule
\textbf{Operación} & \texttt{enviarRespuesta(mensaje)} \\
\textbf{Referencias cruzadas} & Realiza CU-04 (sincronización por
BroadcastChannel); RF-01.7, RF-02.1. \\
\textbf{Precondiciones} & La Estación de Recepción está abierta; el
personal de recepción o enfermería redactó un texto no vacío. \\
\textbf{Postcondiciones} & El mensaje del personal quedó agregado a
\texttt{chatHistory} en ambas estaciones (Recepción y Paciente), de forma
sincronizada a través del canal \texttt{nexora-chat}, y visible para el
paciente en su propia pantalla. \\
\bottomrule
\end{tabularx}
\caption{Contrato de la operación SO-04 --- enviarRespuesta(mensaje)}
\label{tab:col-contrato-so04}
\end{table}

\subsubsection{Diagrama de colaboración DC-04}

\diagrama{08_colaboracion_dc04}{0.85\textwidth}{DC-04 --- Colaboración para enviarRespuesta(mensaje)}{fig:col-dc04}

\subsubsection{Patrones GRASP aplicados}

\begin{itemize}[leftmargin=*, itemsep=2pt]
    \item \textbf{Controlador:} \texttt{RecepcionPage} recibe el evento de
    envío desde \texttt{ReceptionControls} y lo delega a
    \texttt{useChatChannel}, sin conocer los detalles de sincronización
    entre estaciones.
    \item \textbf{Indirección:} la \texttt{BroadcastChannel} nativa del
    navegador media entre las dos instancias de \texttt{useChatChannel}
    (Recepción y Paciente), que nunca se referencian directamente entre sí.
    \item \textbf{Bajo acoplamiento:} \texttt{ReceptionControls} solo conoce
    un \textit{callback} \texttt{onSend(texto)}; no sabe nada sobre el canal
    de sincronización ni sobre la estación del paciente.
    \item \textbf{Alta cohesión:} \texttt{useChatChannel} concentra toda la
    responsabilidad de publicar y recibir mensajes por el canal, tanto del
    lado de Recepción como del lado del Paciente (misma implementación,
    instanciada en cada estación).
\end{itemize}

\subsubsection{Objetos del diagrama y su ubicación en el código}

\begin{table}[H]
\centering
\begin{tabularx}{\textwidth}{@{}l X@{}}
\toprule
\textbf{Objeto} & \textbf{Ubicación real} \\
\midrule
\texttt{ReceptionControls} & Componente $\rightarrow$ \texttt{web/components/reception/ReceptionControls.tsx} \\
\texttt{RecepcionPage} & Componente $\rightarrow$ \texttt{web/app/recepcion/page.tsx} \\
\texttt{useChatChannel (Recepción)} & Hook $\rightarrow$ \texttt{web/hooks/useChatChannel.ts} \\
\texttt{BroadcastChannel} & API nativa del navegador (canal \texttt{nexora-chat}) \\
\texttt{useChatChannel (Paciente)} & Hook $\rightarrow$ \texttt{web/hooks/useChatChannel.ts} \\
\texttt{ChatBubble (Paciente)} & Componente $\rightarrow$ \texttt{web/components/chat/ChatBubble.tsx} \\
\bottomrule
\end{tabularx}
\caption{Trazabilidad de objetos --- DC-04}
\label{tab:col-objetos-dc04}
\end{table}

\subsection{SO-05: capturarMuestra(frames, etiqueta) --- DC-05}

\subsubsection{Contrato de operación}

\begin{table}[H]
\centering
\begin{tabularx}{\textwidth}{@{}p{4.4cm} X@{}}
\toprule
\textbf{Operación} & \texttt{capturarMuestra(frames, etiqueta)} \\
\textbf{Referencias cruzadas} & Realiza CU-06 (grabación de una repetición);
RF-T-03, RF-T-04. \\
\textbf{Precondiciones} & La Estación del Paciente está en Modo Captura; se
seleccionó o creó una etiqueta de seña; la cuenta regresiva de la repetición
ya finalizó. \\
\textbf{Postcondiciones} & Si los datos son válidos (etiqueta bien formada,
cantidad de \textit{frames} y de \textit{blendshapes} consistente con el
modo activo), la muestra quedó persistida como archivo \texttt{.npy} en el
árbol de datos correspondiente al modo (legado o Tasks), asociada a la
etiqueta, y el conteo de muestras por etiqueta quedó actualizado. Si los
datos no son válidos, no se guardó ningún archivo y el error quedó
informado al Modo Captura. \\
\bottomrule
\end{tabularx}
\caption{Contrato de la operación SO-05 --- capturarMuestra(frames, etiqueta)}
\label{tab:col-contrato-so05}
\end{table}

\subsubsection{Diagrama de colaboración DC-05}

\diagrama{08_colaboracion_dc05}{\textwidth}{DC-05 --- Colaboración para capturarMuestra(frames, etiqueta)}{fig:col-dc05}

\subsubsection{Patrones GRASP aplicados}

\begin{itemize}[leftmargin=*, itemsep=2pt]
    \item \textbf{Controlador:} \texttt{CaptureMode} coordina la secuencia de
    grabación (cuenta regresiva, acumulación de \textit{frames}, envío) como
    punto único de entrada de la operación en el frontend.
    \item \textbf{Experto en información:} \texttt{capture\_validation}
    concentra el conocimiento de las reglas de validación (formato de
    etiqueta, consistencia de \textit{blendshapes} según el modo) y es quien
    debe decidir si los datos son aceptables, no el \textit{handler}
    \texttt{capture}.
    \item \textbf{Creador:} el \textit{handler} \texttt{capture} es
    responsable de crear el archivo \texttt{.npy} persistido, porque es
    quien contiene los datos ya validados y quien conoce la ruta de destino
    según el modo (legado o Tasks).
    \item \textbf{Bajo acoplamiento:} \texttt{CaptureMode} no conoce las
    reglas de validación ni el formato de persistencia en disco; solo conoce
    el contrato HTTP de \texttt{/capture} y su respuesta.
\end{itemize}

\subsubsection{Objetos del diagrama y su ubicación en el código}

\begin{table}[H]
\centering
\begin{tabularx}{\textwidth}{@{}l X@{}}
\toprule
\textbf{Objeto} & \textbf{Ubicación real} \\
\midrule
\texttt{CaptureMode} & Componente $\rightarrow$ \texttt{web/components/capture/CaptureMode.tsx} \\
\texttt{captureHandlerRef} & Ref (manejador de \textit{frames} en modo captura) $\rightarrow$ \texttt{web/components/capture/CaptureMode.tsx} \\
\texttt{capture} & Handler FastAPI $\rightarrow$ \texttt{app.py} \\
\texttt{capture\_validation} & Módulo de funciones de validación $\rightarrow$ \texttt{capture\_validation.py} \\
\texttt{dataset/custom/<etiqueta>/*.npy} & Archivo persistido en disco $\rightarrow$ \texttt{dataset/custom/} (o \texttt{dataset/custom\_tasks/} + \texttt{dataset/custom\_blendshapes/} en Modo Tasks) \\
\bottomrule
\end{tabularx}
\caption{Trazabilidad de objetos --- DC-05}
\label{tab:col-objetos-dc05}
\end{table}

\subsection{SO-06: otorgarConsentimiento(paciente) --- DC-06}
\textit{Diseñado --- pendiente de implementación. Cubre el requisito oficial RF-03.}

\subsubsection{Contrato de operación}

\begin{table}[H]
\centering
\begin{tabularx}{\textwidth}{@{}p{4.4cm} X@{}}
\toprule
\textbf{Operación} & \texttt{otorgarConsentimiento(paciente)} \\
\textbf{Referencias cruzadas} & Realiza CU-09 completo; RF-03.1. \\
\textbf{Precondiciones} & La Estación del Paciente muestra la Pantalla
P-05; la cámara todavía no fue activada. \\
\textbf{Postcondiciones} & El consentimiento quedó registrado (sin datos
identificatorios) en el registro de auditoría, y \texttt{ConsentGate}
autorizó el montaje de los \textit{hooks} de captura (\texttt{useMediaPipe}
/ \texttt{useMediaPipeTasks}), habilitando la Pantalla P-02. \\
\bottomrule
\end{tabularx}
\caption{Contrato de la operación SO-06 --- otorgarConsentimiento(paciente)}
\label{tab:col-contrato-so06}
\end{table}

\subsubsection{Diagrama de colaboración DC-06}

\diagrama{08_colaboracion_dc06}{0.9\textwidth}{DC-06 --- Colaboración para otorgarConsentimiento(paciente) --- diseñado, pendiente de implementación}{fig:col-dc06}

\subsubsection{Patrones GRASP aplicados}

\begin{itemize}[leftmargin=*, itemsep=2pt]
    \item \textbf{Controlador:} \texttt{ConsentGate} recibe la confirmación
    desde la Pantalla de Consentimiento y coordina el registro del evento
    y la habilitación de la cámara, como único punto de entrada de la
    operación en el frontend.
    \item \textbf{Experto en información:} \texttt{audit\_log} es quien
    conoce el formato de persistencia de los eventos anónimos, por lo que
    es el responsable de escribirlos, en vez de que cada componente sepa
    escribir su propio archivo.
    \item \textbf{Indirección:} \texttt{audit\_log} media entre el
    \textit{handler} \texttt{consent} y el archivo físico
    \texttt{logs/consent.jsonl}, igual que \texttt{capture\_validation}
    media entre \texttt{capture} y el sistema de archivos en DC-05.
    \item \textbf{Bajo acoplamiento:} la Pantalla de Consentimiento no
    conoce el contrato HTTP ni el formato de persistencia; solo invoca a
    \texttt{ConsentGate} con la decisión del paciente.
\end{itemize}

\subsubsection{Objetos del diagrama y su ubicación en el código}

\begin{table}[H]
\centering
\begin{tabularx}{\textwidth}{@{}l X@{}}
\toprule
\textbf{Objeto} & \textbf{Ubicación real (diseñada)} \\
\midrule
\texttt{PantallaConsentimiento (P-05)} & Componente nuevo $\rightarrow$ previsto en \texttt{web/components/consent/} \\
\texttt{ConsentGate} & Componente de guardia nuevo $\rightarrow$ previsto antepuesto a \texttt{web/app/paciente/page.tsx} \\
\texttt{consent} & Handler FastAPI nuevo $\rightarrow$ previsto en \texttt{app.py} (\texttt{POST /consent}) \\
\texttt{audit\_log} & Módulo de persistencia de eventos nuevo $\rightarrow$ previsto en \texttt{audit\_log.py} \\
\texttt{logs/consent.jsonl} & Archivo persistido en disco $\rightarrow$ previsto en \texttt{logs/} \\
\bottomrule
\end{tabularx}
\caption{Trazabilidad de objetos --- DC-06 (diseñado, pendiente de implementación)}
\label{tab:col-objetos-dc06}
\end{table}

\subsection{SO-07: registrarEventoAuditoria(operacion) --- DC-07}
\textit{Diseñado --- pendiente de implementación. Cubre el requisito oficial RF-04.}

\subsubsection{Contrato de operación}

\begin{table}[H]
\centering
\begin{tabularx}{\textwidth}{@{}p{4.4cm} X@{}}
\toprule
\textbf{Operación} & \texttt{registrarEventoAuditoria(operacion)} \\
\textbf{Referencias cruzadas} & Realiza CU-11, incluido por CU-01 y CU-04; RF-04.1. \\
\textbf{Precondiciones} & Un \textit{handler} del backend (\texttt{websocket\_predict},
\texttt{process\_sentence} o \texttt{capture}) completó una operación de
sistema. \\
\textbf{Postcondiciones} & Un evento anónimo (tipo de operación, marca
temporal, resultado) quedó persistido en el registro de auditoría, sin
datos identificatorios ni imágenes o video. \\
\bottomrule
\end{tabularx}
\caption{Contrato de la operación SO-07 --- registrarEventoAuditoria(operacion)}
\label{tab:col-contrato-so07}
\end{table}

\subsubsection{Diagrama de colaboración DC-07}

\diagrama{08_colaboracion_dc07}{0.9\textwidth}{DC-07 --- Colaboración para registrarEventoAuditoria(operacion) --- diseñado, pendiente de implementación}{fig:col-dc07}

\subsubsection{Patrones GRASP aplicados}

\begin{itemize}[leftmargin=*, itemsep=2pt]
    \item \textbf{Experto en información:} \texttt{audit\_log} concentra el
    conocimiento del formato del evento anónimo, por lo que es el único
    responsable de armarlo y persistirlo.
    \item \textbf{Bajo acoplamiento:} ningún \textit{handler} existente
    (\texttt{websocket\_predict}, \texttt{process\_sentence},
    \texttt{capture}) necesita conocer el formato de persistencia de
    auditoría: cada uno invoca a \texttt{audit\_log} con el tipo de
    operación y el resultado, igual para los tres.
    \item \textbf{Alta cohesión:} toda la responsabilidad de anonimizar,
    formatear y persistir el evento queda concentrada en
    \texttt{audit\_log}, sin dispersarse entre los \textit{handlers}.
\end{itemize}

\subsubsection{Objetos del diagrama y su ubicación en el código}

\begin{table}[H]
\centering
\begin{tabularx}{\textwidth}{@{}l X@{}}
\toprule
\textbf{Objeto} & \textbf{Ubicación real (diseñada)} \\
\midrule
\texttt{handler (websocket\_predict / process\_sentence / capture)} & Handlers FastAPI existentes $\rightarrow$ \texttt{app.py} (cada uno invocaría a \texttt{audit\_log} al finalizar) \\
\texttt{audit\_log} & Módulo de persistencia de eventos nuevo $\rightarrow$ previsto en \texttt{audit\_log.py} \\
\texttt{logs/audit.jsonl} & Archivo persistido en disco $\rightarrow$ previsto en \texttt{logs/} \\
\bottomrule
\end{tabularx}
\caption{Trazabilidad de objetos --- DC-07 (diseñado, pendiente de implementación)}
\label{tab:col-objetos-dc07}
\end{table}

\subsection{SO-08: reportarTraduccion(mensaje, motivo) --- DC-08}
\textit{Diseñado --- pendiente de implementación. Cubre el requisito oficial RF-06.}

\subsubsection{Contrato de operación}

\begin{table}[H]
\centering
\begin{tabularx}{\textwidth}{@{}p{4.4cm} X@{}}
\toprule
\textbf{Operación} & \texttt{reportarTraduccion(mensaje, motivo)} \\
\textbf{Referencias cruzadas} & Realiza CU-10, extiende a CU-01 y CU-04; RF-06.1. \\
\textbf{Precondiciones} & Existe un mensaje visible en el panel de
conversación de P-02 o P-03; el actor hizo clic en su ícono de reporte. \\
\textbf{Postcondiciones} & El reporte (mensaje señalado y motivo) quedó
persistido, sin datos identificatorios del paciente, y el actor recibió
confirmación. \\
\bottomrule
\end{tabularx}
\caption{Contrato de la operación SO-08 --- reportarTraduccion(mensaje, motivo)}
\label{tab:col-contrato-so08}
\end{table}

\subsubsection{Diagrama de colaboración DC-08}

\diagrama{08_colaboracion_dc08}{0.9\textwidth}{DC-08 --- Colaboración para reportarTraduccion(mensaje, motivo) --- diseñado, pendiente de implementación}{fig:col-dc08}

\subsubsection{Patrones GRASP aplicados}

\begin{itemize}[leftmargin=*, itemsep=2pt]
    \item \textbf{Controlador:} la página de estación (\texttt{PacientePage}
    o \texttt{RecepcionPage}) recibe el evento de reporte desde
    \texttt{ChatBubble} y lo delega al backend, sin conocer el formato de
    persistencia.
    \item \textbf{Creador:} el \textit{handler} \texttt{report} es
    responsable de crear el registro persistido, porque es quien recibe el
    mensaje señalado y el motivo ya validados.
    \item \textbf{Bajo acoplamiento:} \texttt{ChatBubble} solo conoce un
    \textit{callback} \texttt{onReport(mensaje, motivo)}; no sabe nada del
    contrato HTTP ni de la persistencia.
    \item \textbf{Indirección:} \texttt{audit\_log} media entre el
    \textit{handler} \texttt{report} y el archivo
    \texttt{logs/reports.jsonl}, reutilizando el mismo módulo de
    persistencia que SO-06 y SO-07.
\end{itemize}

\subsubsection{Objetos del diagrama y su ubicación en el código}

\begin{table}[H]
\centering
\begin{tabularx}{\textwidth}{@{}l X@{}}
\toprule
\textbf{Objeto} & \textbf{Ubicación real (diseñada)} \\
\midrule
\texttt{ChatBubble} & Componente existente, con ícono de reporte nuevo $\rightarrow$ \texttt{web/components/chat/ChatBubble.tsx} \\
\texttt{PacientePage / RecepcionPage} & Componentes existentes $\rightarrow$ \texttt{web/app/paciente/page.tsx}, \texttt{web/app/recepcion/page.tsx} \\
\texttt{report} & Handler FastAPI nuevo $\rightarrow$ previsto en \texttt{app.py} (\texttt{POST /report}) \\
\texttt{audit\_log} & Módulo de persistencia de eventos nuevo $\rightarrow$ previsto en \texttt{audit\_log.py} \\
\texttt{logs/reports.jsonl} & Archivo persistido en disco $\rightarrow$ previsto en \texttt{logs/} \\
\bottomrule
\end{tabularx}
\caption{Trazabilidad de objetos --- DC-08 (diseñado, pendiente de implementación)}
\label{tab:col-objetos-dc08}
\end{table}
